\section*{अध्याय \ref{ch:g11:seq} --- अनुक्रम: पहला पाठ्यक्रम}

\begin{solution}{exo:g11:seq:1}
$u_n = \frac{n}{n+1}$: $u_1 = \frac12$, $u_2 = \frac23$, $u_3 = \frac34$।

$u_0 = 5$, $u_{n+1} = 3u_n - 2$: $u_1 = 13$, $u_2 = 37$, $u_3 = 109$।

$u_n = (-1)^n n$: $u_1 = -1$, $u_2 = 2$, $u_3 = -3$।
\end{solution}

\begin{solution}{exo:g11:seq:2}
$u_{10} = 7 + 10 \times (-3) = -23$ और $u_{25} = 7 - 75 = -68$।

$(v_n)$ के लिए: $v_8 = v_3 + 5d$ से $26 = 11 + 5d$ मिलता है, इसलिए $d = 3$; फिर $v_0 = v_3 - 3d = 11 - 9 = 2$।
\end{solution}

\begin{solution}{exo:g11:seq:3}
$u_8 = 5 \times 2^8 = 1280$।

$(v_n)$ के लिए: $v_4 = v_2\, q^2$ से $48 = 12 q^2$ मिलता है, इसलिए $q^2 = 4$ और $q = 2$ (पद धनात्मक हैं)।
फिर $v_0 = \frac{v_2}{q^2} = \frac{12}{4} = 3$।
\end{solution}

\begin{solution}{exo:g11:seq:4}
$1 + \dots + 500 = \frac{500 \times 501}{2} = 125\,250$।

$4 + 7 + \dots + 61$ $d = 3$ और $\frac{61 - 4}{3} + 1 = 20$ पदों वाला \omterm{def:g11:seq:arithmetic}{समांतर अनुक्रम} है: योग $20 \times \frac{4 + 61}{2} = 650$।

$1 + \frac12 + \dots + \frac{1}{2^{10}}$ $q = \frac12$ और $11$ पदों वाला \omterm{def:g11:seq:geometric}{गुणोत्तर अनुक्रम} है: $\frac{1 - (1/2)^{11}}{1 - 1/2} = 2\left(1 - \frac{1}{2048}\right)
= \frac{2047}{1024}$।
\end{solution}

\begin{solution}{exo:g11:seq:5}
$u_{n+1} - u_n = 4(n+1) - 1 - 4n + 1 = 4$: $d = 4$ वाला \omterm{def:g11:seq:arithmetic}{समांतर}।

$\frac{v_{n+1}}{v_n} = \frac{2^{n+1}}{3^{n+2}} \cdot \frac{3^{n+1}}{2^n}
= \frac23$: $q = \frac23$ वाला \omterm{def:g11:seq:geometric}{गुणोत्तर}।

$w_0 = 0$, $w_1 = 2$, $w_2 = 6$: अंतर $2$ और $4$ अलग-अलग हैं, इसलिए \omterm{def:g11:seq:arithmetic}{समांतर} नहीं; $\frac{w_1}{w_0}$ तो
परिभाषित ही नहीं है, और अनुपात $\frac{w_2}{w_1} = 3 \neq \frac{w_3}{w_2} = 2$: इनमें से कोई नहीं।
\end{solution}

\begin{solution}{exo:g11:seq:6}
पंक्तियों के आकार \omterm{def:g11:seq:arithmetic}{समांतर} हैं: पहला पद $16$, अंतर $2$। अंतिम (बीसवीं) पंक्ति
में $16 + 19 \times 2 = 54$ सीटें हैं। कुल $20 \times \frac{16 + 54}{2} = 700$ सीटें हैं।
\end{solution}

\begin{solution}{exo:g11:seq:7}
$n$ घंटों बाद समष्टि $500 \times 2^n$ है। रात 8 बजे $n = 8$: $500 \times 256 = 128\,000$ जीवाणु। हमें $500 \times 2^n >
10^6$ चाहिए,
अर्थात् $2^n > 2000$: चूँकि $2^{10} = 1024$ और $2^{11} = 2048$ हैं, इसलिए समष्टि पहली बार $11$ पूरे घंटों
के बाद, यानी रात 11 बजे, दस लाख से आगे निकलती है।
\end{solution}

\begin{solution}{exo:g11:seq:8}
$c_2 = 1.002 \times 100 + 100 = 200.20$ और $c_3 = 1.002 \times 200.20 + 100 \approx 300.60$। अंतर $c_2 - c_1 = 100.20$ और $c_3 - c_2 \approx 100.40$ बराबर नहीं हैं, इसलिए $(c_n)$ \omterm{def:g11:seq:arithmetic}{समांतर} नहीं है;
अनुपात $\frac{c_2}{c_1} = 2.002$ और $\frac{c_3}{c_2} \approx 1.50$ भी बराबर नहीं हैं, इसलिए वह \omterm{def:g11:seq:geometric}{गुणोत्तर} भी नहीं है। (इस
जैसी मिली-जुली ``गुणा करो फिर जोड़ो'' वाली पुनरावृत्तियाँ \cref{exo:g11:seq:11} की
सहायक-अनुक्रम वाली तरकीब से हल होती हैं।)
\end{solution}

\begin{solution}{exo:g11:seq:9}
$u_{n+1} - u_n = (n+1)^2 - 8(n+1) - n^2 + 8n = 2n - 7$: $n \leq 3$ के लिए ऋणात्मक, $n \geq 4$ के लिए धनात्मक। इसलिए $(u_n)$ घटकर $u_4 = 16 - 32 = -16$ तक
जाता है, फिर बढ़ता है: वह एकदिष्ट नहीं है।

$(v_n)$ के पद धनात्मक हैं और
\[
\frac{v_{n+1}}{v_n} = \frac{3^{n+1}}{(n+1)!} \cdot \frac{n!}{3^n}
= \frac{3}{n+1},
\]
जो $n \leq 1$ के लिए $> 1$, $n = 2$ के लिए $= 1$, और $n \geq 3$ के लिए $< 1$ है: \omterm{def:g11:seq:sequence}{अनुक्रम} $v_2 = v_3 = \frac92$
तक बढ़ता है, फिर घटता है।
\end{solution}

\begin{solution}{exo:g11:seq:10}
पहले $n$ पद $u_0, \dots, u_{n-1}$ हैं, जहाँ $u_0 = 3$ और $u_{n-1} = 3 + 4(n-1) = 4n - 1$। उनका योग है
\[
n \times \frac{3 + (4n-1)}{2} = n(2n + 1) = 903,
\]
इसलिए $2n^2 + n - 903 = 0$। यहाँ $\Delta = 1 + 4 \times 2 \times 903 =
7225 = 85^2$, और $n = \frac{-1 + 85}{4} = 21$ (ऋणात्मक \omterm{def:g11:quad:discriminant}{मूल} छोड़ दिया गया)। जाँच: $21 \times 43 = 903$।
\end{solution}

\begin{solution}{exo:g11:seq:11}
\emph{1.} $u_1 = 5$, $u_2 = 9$, $u_3 = 17$: हर पद $4, 8, 16$ से एक अधिक है, जिससे $u_n = 2^{n+1} + 1$ का
अनुमान बनता है।

\emph{2.} $v_n = u_n - 1$ के साथ:
\[
v_{n+1} = u_{n+1} - 1 = 2u_n - 1 - 1 = 2(u_n - 1) = 2v_n,
\]
इसलिए $(v_n)$ अनुपात $2$ और पहला पद $v_0 = u_0 - 1 = 2$ वाला \omterm{def:g11:seq:geometric}{गुणोत्तर अनुक्रम} है।

\emph{3.} अतः $v_n = 2 \times 2^n = 2^{n+1}$ और $u_n = v_n + 1 = 2^{n+1} + 1$, जिससे अनुमान की पुष्टि हो जाती है। ($v_n$ में
घटाई गई संख्या $1$ $x \mapsto 2x - 1$ का \omterm{pb:g10:functions:1}{स्थिर बिंदु} है; यही विचार \cref{ch:g12:seq} में $u_{n+1} = au_n + b$ के
लिए फिर लौटता है।)
\end{solution}

\begin{solution}{pb:g11:seq:1}
\textbf{1.} $h_1 = 1$, $h_2 = 3$, $h_3 = 7$।

\textbf{2.} सबसे बड़ी तश्तरी हटाने के लिए उसके ऊपर की $n$ तश्तरियों को
पहले बची हुई खूँटी पर जाना पड़ता है ($h_n$ चालें); बड़ी तश्तरी पार होती है
($1$ चाल); फिर $n$ तश्तरियों को उसके ऊपर वापस चढ़ना पड़ता है ($h_n$
चालें): $h_{n+1} = 2h_n + 1$।

\textbf{3.} $v_{n+1} = h_{n+1} + 1 = 2h_n + 2 = 2v_n$: $v_1 = 2$ के साथ अनुपात $2$ वाला \omterm{def:g11:seq:geometric}{गुणोत्तर}, इसलिए $v_n = 2^n$ और
$h_n = 2^n - 1$।

\textbf{4.} $2^{64} - 1 \approx 1.8 \times 10^{19}$ सेकंड; प्रति वर्ष $3 \times 10^7$ सेकंड से भाग देने पर: लगभग $6 \times 10^{11}$
वर्ष --- छह सौ अरब वर्ष, यानी ब्रह्मांड की आयु का चालीस गुना। भिक्षु चाय के
विश्राम ले सकते हैं।

\textbf{5.} किसी भी वैध हल में सबसे नीचे वाली तश्तरी की पहली चाल पर विचार
कीजिए: उस क्षण बाक़ी $n$ तश्तरियों का बची हुई एक ही खूँटी पर होना ज़रूरी है
(उन्हें वहाँ पहुँचाने में कम से कम $h_n$ चालें), और नीचे वाली तश्तरी की
अंतिम चाल के बाद उन सबको उसके ऊपर लौटना पड़ता है (कम से कम $h_n$ चालें और):
किसी भी हल को कम से कम $2h_n + 1$ चालें चाहिए। पुनरावृत्ति छत ही नहीं, फ़र्श भी
है: $2^n - 1$ इष्टतम है।

\textbf{6.} $1, 1, 2, 3, 5, 8, 13, 21, 34, 55, 89, 144$।

\textbf{7.} \omterm{def:g11:seq:arithmetic}{समांतर} नहीं ($2 - 1 = 1$, पर $3 - 2 = 1$, $5 - 3 = 2$: अंतर बदलते हैं); \omterm{def:g11:seq:geometric}{गुणोत्तर}
नहीं ($\frac21 = 2$, पर $\frac32 = 1.5$)। \omterm{def:g11:func:monotone}{वर्धमान}: $n \geq 2$ के लिए $F_{n+1} - F_n = F_{n-1} > 0$।

\textbf{8.} $F_k = F_{k+2} - F_{k+1}$, इसलिए
\[
\sum_{k=1}^{n} F_k = (F_3 - F_2) + (F_4 - F_3) + \dots +
(F_{n+2} - F_{n+1}) = F_{n+2} - F_2 = F_{n+2} - 1 .
\]
$n = 6$ के लिए: $1 + 1 + 2 + 3 + 5 + 8 = 20 = F_8 - 1 = 21 - 1$।

\textbf{9.} $F_k F_{k+1} - F_{k-1} F_k = F_k (F_{k+1} -
F_{k-1}) = F_k \cdot F_k = F_k^2$; जोड़ने पर क्रमिक निरसन से $F_n F_{n+1} - F_1 F_0$ मिलता है ($F_0 = 0$ के साथ):
वर्गों का योग $F_n F_{n+1}$ है। $n = 4$ के लिए: $1 + 1 + 4 + 9 = 15 = F_4 F_5 =
3 \times 5$।

\textbf{10.} $F_5 F_3 - F_4^2 = 5 \times 2 - 9 = 1$; $F_6 F_4 - F_5^2 = 8 \times 3 - 25 = -1$; $F_7 F_5 - F_6^2 = 13 \times 5 - 64 = 1$: बारी-बारी से $\pm 1$। $F_{n+1} F_{n-1}$ और $F_n^2$ के बीच
का यही एक इकाई का अंतर वह जादूगर वाला खोया-या-पाया गया वर्ग-इकाई है: $F_n \times F_n$
के वर्ग को काटकर $F_{n+1} \times F_{n-1}$ के आयत में जोड़ने पर एक इकाई बननी या निगली जानी ही
है --- वही पतली फाँक।

\textbf{11.} $F_{n+2} = F_{n+1} + F_n \geq F_n + F_n = 2F_n$ (\omterm{def:g11:seq:sequence}{अनुक्रम} बढ़ता है): हर दो सूचकों में कम से कम दुगुना ---
यानी प्रति सूचक अनुपात $\sqrt2$ वाले \omterm{def:g11:seq:geometric}{गुणोत्तर} से कम नहीं।

\textbf{12.} $1.5$; $1.667$; $1.6$; $1.625$; $1.615$; $1.619$; $1.618$; $1.618$। यदि
$r_n \to L$ हो: $F_{n+2} = F_{n+1} + F_n$ से, $F_{n+1}$ से भाग देने पर, $r_{n+1} = 1 + \frac{1}{r_n}$, इसलिए $L = 1 + \frac1L$, अर्थात्
$L^2 = L + 1$: $L = \varphi = \frac{1 + \sqrt5}{2}$, जो \cref{pb:g10:algebra:1} का स्वर्ण अनुपात है। ख़रगोश सोने में बढ़ते हैं।

\textbf{13.} $v_{n+1} = u_{n+1} - \ell = a u_n + b - \ell$; चूँकि $\ell = a\ell + b$ है, यह $a(u_n - \ell) = a v_n$ है: अनुपात $a$ वाला
\omterm{def:g11:seq:geometric}{गुणोत्तर}। अतः $v_n = a^n v_0$ और $u_n = a^n (u_0 - \ell) + \ell$।

\textbf{14.} \omterm{pb:g10:functions:1}{स्थिर बिंदु}: $\ell = 1.01\ell - 300$ से $\ell = 30\,000$ मिलता है। इसलिए $d_n = 1.01^n (10\,000 - 30\,000) + 30\,000
= 30\,000 - 20\,000 \times 1.01^n$।

\textbf{15.} $d_n \leq 0$ के लिए $1.01^n \geq 1.5$ चाहिए: $1.01^{40} \approx 1.489$, $1.01^{41} \approx 1.504$: $41$-वीं अदायगी ऋण
चुका देती है (और वह $300$ से कुछ कम है)। कुल चुकाई गई राशि: $41 \times 300 = 12\,300$ यूरो से
कुछ ही कम --- उधार लिए गए $10\,000$ पर लगभग $2\,300$ यूरो ब्याज पड़ा।

\textbf{16.} \omterm{pb:g10:functions:1}{स्थिर बिंदु} $\ell = \frac{1000}{1 - 1.02} =
-50\,000$, इसलिए $p_n = 1.02^n \times 100\,000 - 50\,000$। $10$ वर्षों बाद: $1.02^{10} \approx 1.219$: $p_{10} \approx 71\,900$
निवासी।

\textbf{17.} $\frac{1000 \times 1001}{2} = 500\,500$; और $2^{20} - 1 = 1\,048\,575$।

\textbf{18.} $7$ से $502$ तक $5$ के डग में: $\frac{502 - 7}{5} + 1 = 100$ पद; योग $= 100 \times \frac{7 + 502}{2} = 25\,450$।

\textbf{19.} शेष $= 100 \times \frac{1.005^{60} - 1}
{1.005 - 1} \approx 100 \times \frac{0.3489}{0.005} \approx
6\,977$ यूरो --- जिसमें $6\,000$ जमा किए गए और लगभग $977$ कमाए
गए: \omterm{def:g11:seq:geometric}{गुणोत्तर} योग बैंक की मातृभाषा हैं।

\textbf{20.} स्पष्ट सूत्र ``$u_{1000}$ क्या है'' का उत्तर तुरंत दे देते हैं;
पुनरावृत्तियाँ बताती हैं कि तंत्र असल में कैसे बदलते हैं --- और कला दूसरे को
पहले में बदल देने की है। \omterm{def:g11:seq:arithmetic}{समांतर अनुक्रम} जोड़ते हैं, \omterm{def:g11:seq:geometric}{गुणोत्तर} गुणा करते हैं,
और हर परिवार के पास अपना योग-सूत्र है (गाउस की जोड़ी बनाना; दुगुना करने वाली
तरकीब)। स्थिर-बिंदु-और-सहायक वाली तरकीब हर एकघात पुनरावृत्ति को \omterm{def:g11:seq:geometric}{गुणोत्तर} में
बदल देती है --- ऋण, जनसंख्याएँ और मीनार, सब इसी के आगे झुक गए। फिबोनाच्ची
किसी भी परिवार को नहीं मानता, फिर भी क्रमिक निरसन की सर्वसमिकाओं ने उसके योग
और वर्ग पकड़ लिए; उसका पूरा चित्र (एक यथार्थ सूत्र, वह स्वर्ण सीमा) अभी और
मज़बूत औज़ारों की प्रतीक्षा में है।
\end{solution}
